% Context from chapters-tex/inverse-problems.tex; for mathematical reference, not compilation.
\section{Examples of Inverse Problems}

We now discuss several inverse problems in imaging that can be solved using quadratic regularization or nonlinear TV.

                                                                
\subsection{Deconvolution}
\label{subsec-deconv}

The blurring operator \eqref{eq-deblurring-operator} is diagonal in the Fourier domain, so quadratic regularization can be solved efficiently using fast Fourier transforms under periodic boundary conditions. We refer to~\eqref{eq-convol-example} and the corresponding explanations. TV regularization does not reduce to a single diagonal Fourier solve and requires an iterative optimization algorithm.


                                                                
\subsection{Inpainting}
\label{sec-inpainting-variational}

For the inpainting problem, the operator defined in \eqref{eq-inp-operator} is diagonal in space
\eq{
	 \Phi = \diag_m( \de_{\Om^c}[m] ),
}
and is an orthogonal projector $\Phi^* = \Phi$.

For noiseless data, we enforce the affine constraint $\enscond{f \in \RR^N}{y=\Phi f}$ using the orthogonal projector
\eq{
	\foralls x, \quad P_y(f)(x) = \choice{
		f(x) \qifq x \in \Om,\\
		\; y(x) \quad \qifq x \notin \Om.
	}
}


Projected gradient descent solves~\eqref{eq-ip-regul-noiseless} for a smooth prior.
For the Sobolev energy, the algorithm iterates
\eq{
	\iit{f}= P_y( \it{f} + \tau \Delta \it{f} ).
}
The iteration converges if $0<\tau<1/4$, since $\norm{\Delta}\leq8$. Figure \ref{fig-inp-sob-iter} shows how it progressively fills the missing region.

                                                                                                         

                                        
   
                                                                                                                                                                                                           
                          

\myfigure{
\tabquatre{
\image{invpbm}{.24}{sob-inp-iter-1}&
\image{invpbm}{.24}{sob-inp-iter-2}&
\image{invpbm}{.24}{sob-inp-iter-3}&
\image{invpbm}{.24}{sob-inp-iter-4} \\
$k=1$ & $k=10$ & $k=20$ & $k=100$ 
}
}{ 
	Sobolev projected gradient descent algorithm.   
}{fig-inp-sob-iter}

           
          
                                        
                                            
                                        
                                                
  
    
                                                               
                     

Figure \ref{fig-inpainting-parrot} illustrates the use of Sobolev inpainting to remove a cage from an image.

\myfigure{
\tabtrois{
\image{invpbm}{.3}{inpainting-parrot-image}&
\image{invpbm}{.3}{inpainting-parrot-mask}&
\image{invpbm}{.3}{inpainting-parrot-sobolev}\\
Image $f_0$ & Observation $y=\Phi f_0$ & Sobolev $f^\star$ \\
}
}{ 
    Removing a parrot cage by inpainting.  
}{fig-inpainting-parrot}

For the smoothed TV prior, the gradient descent reads
\eq{
	\iit{f}= P_y\pa{
			\it{f} + \tau \diverg\pa{ \frac{ \nabla \it{f}}{ \sqrt{\epsilon^2+\norm{\nabla \it{f}}^2} } }
		}
}
which converges if $0<\tau<\epsilon/4$.

Figure \ref{fig-inpainting-cameraman} compares Sobolev inpainting with smoothed TV inpainting for small $\epsilon$. In this example, TV gives a modest visual improvement without increasing the SNR.

                                                                                                                                                                

\myfigure{
\tabdeux{
\image{invpbm}{.3}{inpainting-cameraman-image}&
\image{invpbm}{.3}{inpainting-cameraman-observations}\\
Image $f_0$ & Observation $y=\Phi f_0$ \\
\image{invpbm}{.3}{inpainting-cameraman-sobolev}&
\image{invpbm}{.3}{inpainting-cameraman-tv} \\
Sobolev $f^\star$ & TV $f^\star$
}
}{ 
	Inpainting with Sobolev and TV regularization.   
}{fig-inpainting-cameraman}



                                                                
\subsection{Tomography Inversion}

In an idealized two-dimensional tomography model, each projection integrates the unknown image along a line $\De_{t,\th}$ defined by
\eq{
	x \cdot \tau_\theta = x_1 \cos\th + x_2\sin\th = t
}
Here $\tau_\theta=(\cos\theta,\sin\theta)$ is the unit normal to the line.

The resulting Radon transform is
\eq{
	\foralls \th \in [0,\pi), \foralls t \in \RR, \quad
	p_{\th}(t) = \int_{\De_{t,\th}} f(x) \,d s
	 = \iint f(x)\, \de( x \cdot \tau_\theta - t )\, d x 
}
see Figure~\ref{fig-tomo-principle}.

\myfigure{
\image{invpbm}{.4}{tomo-principle}
}{ 
	Principle of tomography acquisition.   
}{fig-tomo-principle}

With $\hat f(\omega)=\int f(x)e^{-\imath x\cdot\omega}\,\mathrm{d}x$, the Fourier slice theorem identifies the one-dimensional Fourier transform of each projection with a radial slice of the two-dimensional Fourier transform
\eql{\label{eq-fourier-slice}
	\foralls \th \in [0,\pi)~,~ 
\foralls \xi \in \RR\quad
		\hat p_\th(\xi) = \hat f( \xi \cos \th, \xi \sin \th ).
}
Changing to polar coordinates in Fourier inversion gives the filtered-backprojection formula
\eq{
	f(x) = 
\frac{1}{2\pi} \int_{0}^\pi (p_\th \star h)(x \cdot \tau_\theta)\, d \th
}
with $\hat h(\xi) = |\xi|$.


\myfigure{
\tabtrois{
                                
\image{invpbm}{.3}{tomo-image}&
\image{invpbm}{.35}{tomo-radon-subsample}&
\image{invpbm}{.3}{tomo-radon-fourier}\\
Image $f$ & Radon sub-sampling & Fourier domain
}
}{ 
	Partial Fourier measurements.  
}{fig-tomo-radon-subsample}


We model acquisition at finitely many equally spaced orientations $\{\th_k=\pi k/K\}_{0\leq k<K}$ by the partial Radon transform
\eq{
	R f = ( p_{\th_k} )_{0 \leq k < K}.
}
By~\eqref{eq-fourier-slice}, knowing $R f$ is equivalent to knowing the Fourier transform of $f$ along the corresponding radial lines:
\eq{
	\{ \hat f(\xi \cos(\th_k), \xi \sin(\th_k)) \}_k.
}
For a simplified discrete model, we therefore treat acquisition as direct sampling of the Fourier transform:
\eq{
	\Phi f = ( \hat f[\om] )_{\om \in \Om} \in \CC^P
}
where $\Om$ consists of discrete radial sampling lines in the Fourier plane; see Figure \ref{fig-tomo-radon-subsample}, right.

In this idealized discrete model, recovery is an inpainting problem in the Fourier domain. We use the unitary discrete Fourier transform and consider measurements
\eq{
	\foralls \om \in \Om, \quad y[\om] = \hat f[\om] + w[\om]
}
where $w[\om]$ denotes measurement noise, modeled here as Gaussian white noise.

\myfigure{
\tabtrois{
\image{invpbm}{.3}{tomo-image}&
\image{invpbm}{.3}{tomo-pinv-13}&
\image{invpbm}{.3}{tomo-pinv-32}\\
Image $f_0$ & 13 projections & 32 